\cotefile{64}
\scanpage{28}
On a

\[
\left\{\begin{aligned}
&M_0\text{ et }\mathbb F_2^{\underline\omega}\text{ sont orthogonaux l'un de l'autre}\\
&\text{dans }\underline{\mathrm{End}}(M_0)\text{ pour la forme }\mathrm{Tr}(uv)
\end{aligned}\right. \tag*{(15)}
\]

\[
\deleted{\unsure{M_0}}\,
\underset{\mathrm{Ker}(\mathrm{End}(M_0)\overset{\mathrm{tr}}{\unsure{\relbar|\to}}\mathbb F_2)}
{\underline{\unsure{\mathfrak{sl}|\mathfrak{gl}}(M_0)}}
=M_0\oplus\mathbb F_2\,\mathrm{id}_{M_0} \tag*{(16)}
\]

\[
\mathbb F_2^{\underline\omega}\cap\mathfrak{sl}(M_0)\simeq
\mathbb F_2\cdot\mathrm{id}_{M_0}\qquad\ldots\qquad\biggr] \tag*{(17)}
\]

Ainsi, la donnée d'un torseur $P_M$ sous $\inserted{V_J(M_0)\simeq}\mathrm{End}(M_0)$ \deleted{\unsure{ou|un}}, équivaut à celle d'un couple de torseurs, \deleted{sous} l'un sous $M_0$, l'autre sous $\mathbb F_2^{\underline\omega}$, soient $I_{\underline M}$ et $S_M$.

\uline{Théorème} Le foncteur \(\displaystyle M\mathrel{\mapstochar\relbar\mapstochar}(\deleted{I}M_0,I_M,S_M)\)

\diagram{\begin{tikzcd}[ampersand replacement=\&,column sep=small]
{(18)} \& {\substack{\mathbb Z/4\mathbb Z\text{-modules}\\ \text{libres de rang }2}}\arrow[r] \& {\text{triplets }(M_0,I_{\deleted{\unsure{M}}},S_{\deleted{\unsure{M}}}).}
\end{tikzcd}}

\deleted{où $M_0$ un vect. de rg $2$} \inserted{avec}

\[
\left\{\begin{aligned}
&M_0\text{ vect. de rg }2\text{ sur }\mathbb F_2\\
&I\quad M_0\text{-torseur}\\
&S\quad\text{torseur sous }\mathbb F_2^{\underline\omega(M_0)}
\end{aligned}\right.
\]

est une équivalence de groupoïdes.

% slanted
\reversemarginpar\margin{\raggedright Équivaut essentiellement à l'isomorphisme $\mathrm{Gl}(2,\mathbb Z/4\mathbb Z)\simeq\overbrace{\mathrm{Gl}(2,\mathbb F_2)}^{\mathfrak S_3}\cdot\mathbb M_2(\mathbb F_2)$ (produit $1/2$ direct), plus l'isomorphisme de $\mathfrak S_3$-modules $(\mathfrak S_3\simeq\mathrm{Gl}(2,\mathbb F_2))$ $\mathbb M_2(\mathbb F_2)\simeq\mathbb F_2^2\times\mathbb F_2^{\underline\omega_0}$, où $\mathrm{Gl}(2,\mathbb F_2)$ opère sur $\mathbb F_2^2$ tautologiquement, et sur $\mathbb F_2^{\underline\omega_0}$ via op. sur $\underline\omega_0$ par signature \ldots\ On a l'inclusion $\mathrm{Sl}(2,\mathbb Z/4\mathbb Z)\unsure{=|\subset|\simeq}\deleted{\unsure{\mathrm{Gl}(2,\mathbb Z/4\mathbb Z)\cdot\mathbb M_2(\mathbb F_2)}}\deleted{\unsure{\mathrm{Gl}(2,\mathbb F_2)\times\mathbb F_2^2\times\mathbb F_2}}\{g\cdot u\in\mathfrak S_3\cdot\mathbb M_2(\mathbb F_2)\mid\mathrm{sg}(g)\,\underline e(\mathrm{Tr}\,u)=1\}$ où $\underline e:\deleted{\unsure{\mathbb Z/2\mathbb Z}}\mathbb F_2\overset{\sim}{\longrightarrow}\{\pm1\}$.}\normalmarginpar

\uline{Corollaire} On a un isom. de suites exactes canoniques ($\forall M$ libre de rg $2$ sur $\Lambda=\mathbb Z/4\mathbb Z$)

\diagram{\begin{tikzcd}[ampersand replacement=\&,column sep=small]
\&1\arrow[r,dash]\&\mathrm{End}_{\mathbb F_2}(M_0)\arrow[r]\&\mathrm{Gl}_{\Lambda}(M)\arrow[r]\&\overset{\simeq\mathfrak S_J}{\mathrm{Gl}(M_0)}\arrow[r,dash]\&1\\
{(19)}\&1\arrow[r,dash]\&M_0\times\mathbb F_2^{\underline\omega}\arrow[u,"\sim"']\arrow[r]\&\mathrm{Aut}(M_0,I_M,S_M)\arrow[u,"\sim"']\arrow[r]\&\mathrm{Gl}(M_0)\arrow[u,equal]\arrow[r,dash]\&1
\end{tikzcd}}

Ici l'hom $\mathrm{End}(M_0)\to\mathrm{Gl}(M)$ est donné par la condition que le composé avec l'hom $\mathrm{End}_{\Lambda}(M)\to\mathrm{End}_{\mathbb F_2}(M_0)$ soit donné par $u\longmapsto\mathrm{id}_M+2u$.
% archive sheet number: 27
